% Pipeline Followed -- Person 2
% Included from report.tex via % Pipeline Followed -- Person 2
% Included from report.tex via % Pipeline Followed -- Person 2
% Included from report.tex via % Pipeline Followed -- Person 2
% Included from report.tex via \input{sections/pipeline_followed}
% Requires in the preamble: booktabs, graphicx.
% Figures referenced: fig:noise_floor, fig:mp_spectrum, fig:threshold_sweep.
% Every number below is written by code into artifacts/ and regenerates with
% `python run_all.py`.

\section{Pipeline Followed}

The network has \textbf{the 60 survey questions as nodes} and \textbf{signed, thresholded
correlations between questions as edges}. A correlation network built from 86 respondents
contains a great deal of structure that is not evidence of anything; every step below
exists to separate the part that is.

\subsection{Parsing, encoding, and missing data}

The file is read with \texttt{encoding='utf-8-sig'} -- a byte-order mark sits on the first
header, and without that flag the column name silently fails to match. Item codes are
parsed from the headers by regular expression rather than by renaming columns, giving four
blocks of exactly 15 items. The Likert scale maps to $-2 \ldots +2$, and the sixth option,
``No Comments'' (37 occurrences), maps to \emph{missing} rather than $0$: it is a refusal
to answer, not the substantive position ``Neutral'', and collapsing the two would place 37
refusals at the exact centre of the scale.

That leaves 542 missing entries (505 blanks plus 37 refusals). How they are distributed
matters more than how many there are:

\begin{center}
\begin{tabular}{lrrrrr}
\toprule
Missing items per respondent & 0 & 1--5 & 6--12 & 13--59 & all 60 \\
\midrule
Respondents & 68 & 17 & 1 & 5 & 5 \\
\bottomrule
\end{tabular}
\end{center}

\noindent
\textbf{495 of the 542 missing cells sit in the 10 respondents with more than 12 missing
items.} Dropping those 10 leaves \textbf{86 respondents with 47 missing cells (0.9\%)},
filled with the column mean. We checked rather than assumed that nothing heavier was
needed: substituting the median moves the correlation standard deviation from 0.1363 to
0.1361 and the edge count at $|r|>0.20$ from 258 to 260, and leaves the number of
significant eigenvalues unchanged. When two simple estimators are indistinguishable in
every downstream result, a regression or Bayesian imputer adds modelling assumptions in
exchange for no measurable change.

\subsection{Row-centring and the association measure}

The grand mean across all responses is $+1.14$ on a $-2 \ldots +2$ scale: respondents
agree with almost everything. Uncorrected, this acquiescence baseline dominates the
similarity structure, because someone who agrees with everything pushes all 60 columns up
together, which is indistinguishable from genuine item association. On the uncentred
matrix, 83\% of the 1{,}770 item correlations are positive and 696 pairs exceed
$|r|>0.20$ -- a 39\% dense graph encoding mostly who is agreeable. We therefore subtract
each respondent's own mean across all 60 items from their row before correlating, leaving
only the \emph{relative} pattern of what they lean toward given their own baseline. After
centring, 45\% of pairs are positive and 258 survive $|r|>0.20$.

Edges are Pearson correlations on that matrix \textbf{with signs retained}; we never take
absolute values. A negative edge says two questions divide the class in opposite
directions, which is a stronger claim than ``these are related'', and retaining signs is
what makes the structural-balance analysis possible at all.

\subsection{The noise floor and spectral filtering}
\label{subsec:mp}

Each of 1{,}770 item correlations is estimated from 86 respondents, so unrelated questions
scatter around zero with standard deviation $1/\sqrt{n-3} = 0.110$; across 1{,}770 pairs
many land high by luck. We quantify this by re-running the whole pipeline on data with the
association destroyed -- every column permuted independently, then centred and correlated
exactly as the real data, 200 times. Observed correlations have SD 0.136 against 0.108 for
the null (Figure~\ref{fig:noise_floor}):

\begin{center}
\begin{tabular}{lrrrr}
\toprule
Threshold $|r| >$ & 0.15 & 0.20 & 0.25 & 0.30 \\
\midrule
Observed edges & 467 & 258 & 111 & 49 \\
Expected by chance & 308 & 119 & 37 & 9 \\
Chance fraction & 66\% & 46\% & 33\% & 19\% \\
\bottomrule
\end{tabular}
\end{center}

PCA alone cannot settle this: PC1 explains only 8.3\% of variance, which distinguishes
weak structure from none not at all. Marchenko--Pastur does. For a correlation matrix from
$T$ observations of $N$ pure-noise variables the eigenvalues fall in an analytic band; at
$T=86$, $N=60$, $Q=1.433$ and $\lambda_{+}=(1+1/\sqrt{Q})^2=3.368$. \textbf{Exactly three
eigenvalues exceed it} -- 4.05, 3.98, 3.53, carrying 19.3\% of variance
(Figure~\ref{fig:mp_spectrum}) -- and the matrix is rebuilt from those three alone with the
diagonal restored to unity. Reading the loadings, they correspond to (i)~systemic
responsibility against technological optimism and traditional instruction,
(ii)~regulatory guardrails against individual autonomy, and (iii)~AI-enabled education and
industry against environmental stewardship. The third is a clean split; the first two are
mixed, and we offer them as readings rather than established constructs. One caveat
belongs with the headline: the fourth eigenvalue is 3.354 against a band edge of 3.368, so
``exactly three'' is correct but holds by 0.4\%.

\subsection{Nodes, edges, and threshold selection}

\textbf{Nodes are the 60 questions; edges are signed correlations above a threshold,
weighted by default and binarised only where a metric requires it} (path lengths, triad
counts). Questions were chosen over respondents deliberately: a respondent network must be
built by a neighbourhood rule such as $k$-nearest-neighbours, and that rule manufactures
its own communities, as Section~\ref{subsec:nulls} shows. An item network is thresholded
directly on a measured quantity and every node carries an interpretable label. Weights are
kept because correlation magnitude is measured information that binarising discards.

The threshold comes from the sweep in Figure~\ref{fig:threshold_sweep}, which rebuilds the
graph at all 46 cutoffs from 0.05 to 0.50. The rule is \emph{the highest cutoff at which
the giant component still spans more than two-thirds of the items and mean degree stays
above 4}, so community and path metrics remain meaningful; since spectral denoising has
already removed the sampling noise the chance table measures, raising the cutoff further
would penalise the same noise twice. This gives $\mathbf{|r|>0.20}$: 124 edges (60
positive, 64 negative), mean degree 4.13, a giant component of 41 of 60 items (68\%), and
modularity 0.405. Two independent controls agree there: the permutation null implies the
raw graph's 258 edges contain about 119 chance edges, leaving roughly 139 real; denoising,
which never sees the null, retains 124 -- an 11\% difference.

Panel~(d) shows why the sweep is necessary rather than decorative. Modularity climbs
steadily to 0.63 as the threshold rises, but by then the graph is 51 components and 47
isolated items. Quoting a single high modularity without the sweep would report
fragmentation as community structure.

\subsection{Null models}
\label{subsec:nulls}

A correlation-derived network has transitivity built into its construction: if A
correlates with B and B with C, A and C are pulled toward correlation automatically. That
alone manufactures clustering and modularity in data with no association whatsoever.
Degree-preserving rewiring and Erd\H{o}s--R\'enyi nulls both take the \emph{finished
graph} as given, so neither can detect structure the construction itself created. Only a
permutation null -- destroying association in the raw responses, then regenerating the
entire chain of centre, correlate, threshold and measure -- tests the claim being made.

This decides the result rather than decorating it. On the respondent network (86 nodes,
286 edges, one component, symmetrised $k$-NN at $k=4$) the observed modularity is 0.416
and average clustering 0.202:

\begin{center}
\begin{tabular}{lrrrl}
\toprule
Null model & Mean & SD & $z$ & Verdict \\
\midrule
Degree-preserving rewiring & 0.329 & 0.011 & 7.74 & passes comfortably \\
Erd\H{o}s--R\'enyi $G(n,m)$ & 0.332 & 0.011 & 7.79 & passes comfortably \\
Column permutation (full pipeline) & 0.400 & 0.019 & 0.79 & \textbf{fails} ($p=0.20$) \\
\bottomrule
\end{tabular}
\end{center}

\noindent
Against the graph-level nulls the community structure looks overwhelming; against the null
that reproduces the construction it disappears. We report this as a finding, not a
failure: those communities are an artefact of correlating and then applying a
neighbourhood rule, which is exactly why the item network is our primary object.

\subsection{Methods tried and rejected}
\label{subsec:rejected}

\begin{center}
\begin{tabular}{p{0.26\linewidth}p{0.40\linewidth}p{0.26\linewidth}}
\toprule
Approach & Measured outcome & Rejected because \\
\midrule
Mutual $k$-NN ($k=4$)
& 58 edges, 38 components, giant component 39.5\%
& Shatters the graph. \\
Raw (uncentred) similarity
& 696 of 1{,}770 pairs above $|r|>0.20$; 39\% density; 83\% positive
& Measures acquiescence, not opinion. \\
Percentile thresholding (p90)
& Items: 177 edges, 1 component. Respondents: 366 edges, 22 components, 20 isolates
& Fixes edge count by construction, so cannot say whether any edge beats chance. \\
Ordinal-regression imputation
& 47.6\% held-out accuracy, MAE 0.66
& Unjustifiable for 47 cells (0.9\%) where mean and median already agree. \\
Keeping 91 respondents rather than 86
& $Q=1.517$, $\lambda_{+}=3.283$, \emph{four} dimensions carrying 25.6\%
& Fourth dimension rests largely on imputed values. \\
\bottomrule
\end{tabular}
\end{center}

\noindent
The last row is a genuine sensitivity result and we report it rather than suppress it: the
``exactly three dimensions'' headline depends on the respondent filter. Retaining the five
respondents who answered between one and 47 items would mean imputing 242 cells with a
model measured at 47.6\% accuracy, and the extra dimension rests disproportionately on
those filled values. The frozen filter is the more defensible choice, but the finding is
conditional on it and a reader should know that.

% Requires in the preamble: booktabs, graphicx.
% Figures referenced: fig:noise_floor, fig:mp_spectrum, fig:threshold_sweep.
% Every number below is written by code into artifacts/ and regenerates with
% `python run_all.py`.

\section{Pipeline Followed}

The network has \textbf{the 60 survey questions as nodes} and \textbf{signed, thresholded
correlations between questions as edges}. A correlation network built from 86 respondents
contains a great deal of structure that is not evidence of anything; every step below
exists to separate the part that is.

\subsection{Parsing, encoding, and missing data}

The file is read with \texttt{encoding='utf-8-sig'} -- a byte-order mark sits on the first
header, and without that flag the column name silently fails to match. Item codes are
parsed from the headers by regular expression rather than by renaming columns, giving four
blocks of exactly 15 items. The Likert scale maps to $-2 \ldots +2$, and the sixth option,
``No Comments'' (37 occurrences), maps to \emph{missing} rather than $0$: it is a refusal
to answer, not the substantive position ``Neutral'', and collapsing the two would place 37
refusals at the exact centre of the scale.

That leaves 542 missing entries (505 blanks plus 37 refusals). How they are distributed
matters more than how many there are:

\begin{center}
\begin{tabular}{lrrrrr}
\toprule
Missing items per respondent & 0 & 1--5 & 6--12 & 13--59 & all 60 \\
\midrule
Respondents & 68 & 17 & 1 & 5 & 5 \\
\bottomrule
\end{tabular}
\end{center}

\noindent
\textbf{495 of the 542 missing cells sit in the 10 respondents with more than 12 missing
items.} Dropping those 10 leaves \textbf{86 respondents with 47 missing cells (0.9\%)},
filled with the column mean. We checked rather than assumed that nothing heavier was
needed: substituting the median moves the correlation standard deviation from 0.1363 to
0.1361 and the edge count at $|r|>0.20$ from 258 to 260, and leaves the number of
significant eigenvalues unchanged. When two simple estimators are indistinguishable in
every downstream result, a regression or Bayesian imputer adds modelling assumptions in
exchange for no measurable change.

\subsection{Row-centring and the association measure}

The grand mean across all responses is $+1.14$ on a $-2 \ldots +2$ scale: respondents
agree with almost everything. Uncorrected, this acquiescence baseline dominates the
similarity structure, because someone who agrees with everything pushes all 60 columns up
together, which is indistinguishable from genuine item association. On the uncentred
matrix, 83\% of the 1{,}770 item correlations are positive and 696 pairs exceed
$|r|>0.20$ -- a 39\% dense graph encoding mostly who is agreeable. We therefore subtract
each respondent's own mean across all 60 items from their row before correlating, leaving
only the \emph{relative} pattern of what they lean toward given their own baseline. After
centring, 45\% of pairs are positive and 258 survive $|r|>0.20$.

Edges are Pearson correlations on that matrix \textbf{with signs retained}; we never take
absolute values. A negative edge says two questions divide the class in opposite
directions, which is a stronger claim than ``these are related'', and retaining signs is
what makes the structural-balance analysis possible at all.

\subsection{The noise floor and spectral filtering}
\label{subsec:mp}

Each of 1{,}770 item correlations is estimated from 86 respondents, so unrelated questions
scatter around zero with standard deviation $1/\sqrt{n-3} = 0.110$; across 1{,}770 pairs
many land high by luck. We quantify this by re-running the whole pipeline on data with the
association destroyed -- every column permuted independently, then centred and correlated
exactly as the real data, 200 times. Observed correlations have SD 0.136 against 0.108 for
the null (Figure~\ref{fig:noise_floor}):

\begin{center}
\begin{tabular}{lrrrr}
\toprule
Threshold $|r| >$ & 0.15 & 0.20 & 0.25 & 0.30 \\
\midrule
Observed edges & 467 & 258 & 111 & 49 \\
Expected by chance & 308 & 119 & 37 & 9 \\
Chance fraction & 66\% & 46\% & 33\% & 19\% \\
\bottomrule
\end{tabular}
\end{center}

PCA alone cannot settle this: PC1 explains only 8.3\% of variance, which distinguishes
weak structure from none not at all. Marchenko--Pastur does. For a correlation matrix from
$T$ observations of $N$ pure-noise variables the eigenvalues fall in an analytic band; at
$T=86$, $N=60$, $Q=1.433$ and $\lambda_{+}=(1+1/\sqrt{Q})^2=3.368$. \textbf{Exactly three
eigenvalues exceed it} -- 4.05, 3.98, 3.53, carrying 19.3\% of variance
(Figure~\ref{fig:mp_spectrum}) -- and the matrix is rebuilt from those three alone with the
diagonal restored to unity. Reading the loadings, they correspond to (i)~systemic
responsibility against technological optimism and traditional instruction,
(ii)~regulatory guardrails against individual autonomy, and (iii)~AI-enabled education and
industry against environmental stewardship. The third is a clean split; the first two are
mixed, and we offer them as readings rather than established constructs. One caveat
belongs with the headline: the fourth eigenvalue is 3.354 against a band edge of 3.368, so
``exactly three'' is correct but holds by 0.4\%.

\subsection{Nodes, edges, and threshold selection}

\textbf{Nodes are the 60 questions; edges are signed correlations above a threshold,
weighted by default and binarised only where a metric requires it} (path lengths, triad
counts). Questions were chosen over respondents deliberately: a respondent network must be
built by a neighbourhood rule such as $k$-nearest-neighbours, and that rule manufactures
its own communities, as Section~\ref{subsec:nulls} shows. An item network is thresholded
directly on a measured quantity and every node carries an interpretable label. Weights are
kept because correlation magnitude is measured information that binarising discards.

The threshold comes from the sweep in Figure~\ref{fig:threshold_sweep}, which rebuilds the
graph at all 46 cutoffs from 0.05 to 0.50. The rule is \emph{the highest cutoff at which
the giant component still spans more than two-thirds of the items and mean degree stays
above 4}, so community and path metrics remain meaningful; since spectral denoising has
already removed the sampling noise the chance table measures, raising the cutoff further
would penalise the same noise twice. This gives $\mathbf{|r|>0.20}$: 124 edges (60
positive, 64 negative), mean degree 4.13, a giant component of 41 of 60 items (68\%), and
modularity 0.405. Two independent controls agree there: the permutation null implies the
raw graph's 258 edges contain about 119 chance edges, leaving roughly 139 real; denoising,
which never sees the null, retains 124 -- an 11\% difference.

Panel~(d) shows why the sweep is necessary rather than decorative. Modularity climbs
steadily to 0.63 as the threshold rises, but by then the graph is 51 components and 47
isolated items. Quoting a single high modularity without the sweep would report
fragmentation as community structure.

\subsection{Null models}
\label{subsec:nulls}

A correlation-derived network has transitivity built into its construction: if A
correlates with B and B with C, A and C are pulled toward correlation automatically. That
alone manufactures clustering and modularity in data with no association whatsoever.
Degree-preserving rewiring and Erd\H{o}s--R\'enyi nulls both take the \emph{finished
graph} as given, so neither can detect structure the construction itself created. Only a
permutation null -- destroying association in the raw responses, then regenerating the
entire chain of centre, correlate, threshold and measure -- tests the claim being made.

This decides the result rather than decorating it. On the respondent network (86 nodes,
286 edges, one component, symmetrised $k$-NN at $k=4$) the observed modularity is 0.416
and average clustering 0.202:

\begin{center}
\begin{tabular}{lrrrl}
\toprule
Null model & Mean & SD & $z$ & Verdict \\
\midrule
Degree-preserving rewiring & 0.329 & 0.011 & 7.74 & passes comfortably \\
Erd\H{o}s--R\'enyi $G(n,m)$ & 0.332 & 0.011 & 7.79 & passes comfortably \\
Column permutation (full pipeline) & 0.400 & 0.019 & 0.79 & \textbf{fails} ($p=0.20$) \\
\bottomrule
\end{tabular}
\end{center}

\noindent
Against the graph-level nulls the community structure looks overwhelming; against the null
that reproduces the construction it disappears. We report this as a finding, not a
failure: those communities are an artefact of correlating and then applying a
neighbourhood rule, which is exactly why the item network is our primary object.

\subsection{Methods tried and rejected}
\label{subsec:rejected}

\begin{center}
\begin{tabular}{p{0.26\linewidth}p{0.40\linewidth}p{0.26\linewidth}}
\toprule
Approach & Measured outcome & Rejected because \\
\midrule
Mutual $k$-NN ($k=4$)
& 58 edges, 38 components, giant component 39.5\%
& Shatters the graph. \\
Raw (uncentred) similarity
& 696 of 1{,}770 pairs above $|r|>0.20$; 39\% density; 83\% positive
& Measures acquiescence, not opinion. \\
Percentile thresholding (p90)
& Items: 177 edges, 1 component. Respondents: 366 edges, 22 components, 20 isolates
& Fixes edge count by construction, so cannot say whether any edge beats chance. \\
Ordinal-regression imputation
& 47.6\% held-out accuracy, MAE 0.66
& Unjustifiable for 47 cells (0.9\%) where mean and median already agree. \\
Keeping 91 respondents rather than 86
& $Q=1.517$, $\lambda_{+}=3.283$, \emph{four} dimensions carrying 25.6\%
& Fourth dimension rests largely on imputed values. \\
\bottomrule
\end{tabular}
\end{center}

\noindent
The last row is a genuine sensitivity result and we report it rather than suppress it: the
``exactly three dimensions'' headline depends on the respondent filter. Retaining the five
respondents who answered between one and 47 items would mean imputing 242 cells with a
model measured at 47.6\% accuracy, and the extra dimension rests disproportionately on
those filled values. The frozen filter is the more defensible choice, but the finding is
conditional on it and a reader should know that.

% Requires in the preamble: booktabs, graphicx.
% Figures referenced: fig:noise_floor, fig:mp_spectrum, fig:threshold_sweep.
% Every number below is written by code into artifacts/ and regenerates with
% `python run_all.py`.

\section{Pipeline Followed}

The network has \textbf{the 60 survey questions as nodes} and \textbf{signed, thresholded
correlations between questions as edges}. A correlation network built from 86 respondents
contains a great deal of structure that is not evidence of anything; every step below
exists to separate the part that is.

\subsection{Parsing, encoding, and missing data}

The file is read with \texttt{encoding='utf-8-sig'} -- a byte-order mark sits on the first
header, and without that flag the column name silently fails to match. Item codes are
parsed from the headers by regular expression rather than by renaming columns, giving four
blocks of exactly 15 items. The Likert scale maps to $-2 \ldots +2$, and the sixth option,
``No Comments'' (37 occurrences), maps to \emph{missing} rather than $0$: it is a refusal
to answer, not the substantive position ``Neutral'', and collapsing the two would place 37
refusals at the exact centre of the scale.

That leaves 542 missing entries (505 blanks plus 37 refusals). How they are distributed
matters more than how many there are:

\begin{center}
\begin{tabular}{lrrrrr}
\toprule
Missing items per respondent & 0 & 1--5 & 6--12 & 13--59 & all 60 \\
\midrule
Respondents & 68 & 17 & 1 & 5 & 5 \\
\bottomrule
\end{tabular}
\end{center}

\noindent
\textbf{495 of the 542 missing cells sit in the 10 respondents with more than 12 missing
items.} Dropping those 10 leaves \textbf{86 respondents with 47 missing cells (0.9\%)},
filled with the column mean. We checked rather than assumed that nothing heavier was
needed: substituting the median moves the correlation standard deviation from 0.1363 to
0.1361 and the edge count at $|r|>0.20$ from 258 to 260, and leaves the number of
significant eigenvalues unchanged. When two simple estimators are indistinguishable in
every downstream result, a regression or Bayesian imputer adds modelling assumptions in
exchange for no measurable change.

\subsection{Row-centring and the association measure}

The grand mean across all responses is $+1.14$ on a $-2 \ldots +2$ scale: respondents
agree with almost everything. Uncorrected, this acquiescence baseline dominates the
similarity structure, because someone who agrees with everything pushes all 60 columns up
together, which is indistinguishable from genuine item association. On the uncentred
matrix, 83\% of the 1{,}770 item correlations are positive and 696 pairs exceed
$|r|>0.20$ -- a 39\% dense graph encoding mostly who is agreeable. We therefore subtract
each respondent's own mean across all 60 items from their row before correlating, leaving
only the \emph{relative} pattern of what they lean toward given their own baseline. After
centring, 45\% of pairs are positive and 258 survive $|r|>0.20$.

Edges are Pearson correlations on that matrix \textbf{with signs retained}; we never take
absolute values. A negative edge says two questions divide the class in opposite
directions, which is a stronger claim than ``these are related'', and retaining signs is
what makes the structural-balance analysis possible at all.

\subsection{The noise floor and spectral filtering}
\label{subsec:mp}

Each of 1{,}770 item correlations is estimated from 86 respondents, so unrelated questions
scatter around zero with standard deviation $1/\sqrt{n-3} = 0.110$; across 1{,}770 pairs
many land high by luck. We quantify this by re-running the whole pipeline on data with the
association destroyed -- every column permuted independently, then centred and correlated
exactly as the real data, 200 times. Observed correlations have SD 0.136 against 0.108 for
the null (Figure~\ref{fig:noise_floor}):

\begin{center}
\begin{tabular}{lrrrr}
\toprule
Threshold $|r| >$ & 0.15 & 0.20 & 0.25 & 0.30 \\
\midrule
Observed edges & 467 & 258 & 111 & 49 \\
Expected by chance & 308 & 119 & 37 & 9 \\
Chance fraction & 66\% & 46\% & 33\% & 19\% \\
\bottomrule
\end{tabular}
\end{center}

PCA alone cannot settle this: PC1 explains only 8.3\% of variance, which distinguishes
weak structure from none not at all. Marchenko--Pastur does. For a correlation matrix from
$T$ observations of $N$ pure-noise variables the eigenvalues fall in an analytic band; at
$T=86$, $N=60$, $Q=1.433$ and $\lambda_{+}=(1+1/\sqrt{Q})^2=3.368$. \textbf{Exactly three
eigenvalues exceed it} -- 4.05, 3.98, 3.53, carrying 19.3\% of variance
(Figure~\ref{fig:mp_spectrum}) -- and the matrix is rebuilt from those three alone with the
diagonal restored to unity. Reading the loadings, they correspond to (i)~systemic
responsibility against technological optimism and traditional instruction,
(ii)~regulatory guardrails against individual autonomy, and (iii)~AI-enabled education and
industry against environmental stewardship. The third is a clean split; the first two are
mixed, and we offer them as readings rather than established constructs. One caveat
belongs with the headline: the fourth eigenvalue is 3.354 against a band edge of 3.368, so
``exactly three'' is correct but holds by 0.4\%.

\subsection{Nodes, edges, and threshold selection}

\textbf{Nodes are the 60 questions; edges are signed correlations above a threshold,
weighted by default and binarised only where a metric requires it} (path lengths, triad
counts). Questions were chosen over respondents deliberately: a respondent network must be
built by a neighbourhood rule such as $k$-nearest-neighbours, and that rule manufactures
its own communities, as Section~\ref{subsec:nulls} shows. An item network is thresholded
directly on a measured quantity and every node carries an interpretable label. Weights are
kept because correlation magnitude is measured information that binarising discards.

The threshold comes from the sweep in Figure~\ref{fig:threshold_sweep}, which rebuilds the
graph at all 46 cutoffs from 0.05 to 0.50. The rule is \emph{the highest cutoff at which
the giant component still spans more than two-thirds of the items and mean degree stays
above 4}, so community and path metrics remain meaningful; since spectral denoising has
already removed the sampling noise the chance table measures, raising the cutoff further
would penalise the same noise twice. This gives $\mathbf{|r|>0.20}$: 124 edges (60
positive, 64 negative), mean degree 4.13, a giant component of 41 of 60 items (68\%), and
modularity 0.405. Two independent controls agree there: the permutation null implies the
raw graph's 258 edges contain about 119 chance edges, leaving roughly 139 real; denoising,
which never sees the null, retains 124 -- an 11\% difference.

Panel~(d) shows why the sweep is necessary rather than decorative. Modularity climbs
steadily to 0.63 as the threshold rises, but by then the graph is 51 components and 47
isolated items. Quoting a single high modularity without the sweep would report
fragmentation as community structure.

\subsection{Null models}
\label{subsec:nulls}

A correlation-derived network has transitivity built into its construction: if A
correlates with B and B with C, A and C are pulled toward correlation automatically. That
alone manufactures clustering and modularity in data with no association whatsoever.
Degree-preserving rewiring and Erd\H{o}s--R\'enyi nulls both take the \emph{finished
graph} as given, so neither can detect structure the construction itself created. Only a
permutation null -- destroying association in the raw responses, then regenerating the
entire chain of centre, correlate, threshold and measure -- tests the claim being made.

This decides the result rather than decorating it. On the respondent network (86 nodes,
286 edges, one component, symmetrised $k$-NN at $k=4$) the observed modularity is 0.416
and average clustering 0.202:

\begin{center}
\begin{tabular}{lrrrl}
\toprule
Null model & Mean & SD & $z$ & Verdict \\
\midrule
Degree-preserving rewiring & 0.329 & 0.011 & 7.74 & passes comfortably \\
Erd\H{o}s--R\'enyi $G(n,m)$ & 0.332 & 0.011 & 7.79 & passes comfortably \\
Column permutation (full pipeline) & 0.400 & 0.019 & 0.79 & \textbf{fails} ($p=0.20$) \\
\bottomrule
\end{tabular}
\end{center}

\noindent
Against the graph-level nulls the community structure looks overwhelming; against the null
that reproduces the construction it disappears. We report this as a finding, not a
failure: those communities are an artefact of correlating and then applying a
neighbourhood rule, which is exactly why the item network is our primary object.

\subsection{Methods tried and rejected}
\label{subsec:rejected}

\begin{center}
\begin{tabular}{p{0.26\linewidth}p{0.40\linewidth}p{0.26\linewidth}}
\toprule
Approach & Measured outcome & Rejected because \\
\midrule
Mutual $k$-NN ($k=4$)
& 58 edges, 38 components, giant component 39.5\%
& Shatters the graph. \\
Raw (uncentred) similarity
& 696 of 1{,}770 pairs above $|r|>0.20$; 39\% density; 83\% positive
& Measures acquiescence, not opinion. \\
Percentile thresholding (p90)
& Items: 177 edges, 1 component. Respondents: 366 edges, 22 components, 20 isolates
& Fixes edge count by construction, so cannot say whether any edge beats chance. \\
Ordinal-regression imputation
& 47.6\% held-out accuracy, MAE 0.66
& Unjustifiable for 47 cells (0.9\%) where mean and median already agree. \\
Keeping 91 respondents rather than 86
& $Q=1.517$, $\lambda_{+}=3.283$, \emph{four} dimensions carrying 25.6\%
& Fourth dimension rests largely on imputed values. \\
\bottomrule
\end{tabular}
\end{center}

\noindent
The last row is a genuine sensitivity result and we report it rather than suppress it: the
``exactly three dimensions'' headline depends on the respondent filter. Retaining the five
respondents who answered between one and 47 items would mean imputing 242 cells with a
model measured at 47.6\% accuracy, and the extra dimension rests disproportionately on
those filled values. The frozen filter is the more defensible choice, but the finding is
conditional on it and a reader should know that.

% Requires in the preamble: booktabs, graphicx.
% Figures referenced: fig:noise_floor, fig:mp_spectrum, fig:threshold_sweep.
% Every number below is written by code into artifacts/ and regenerates with
% `python run_all.py`.

\section{Pipeline Followed}

The network has \textbf{the 60 survey questions as nodes} and \textbf{signed, thresholded
correlations between questions as edges}. A correlation network built from 86 respondents
contains a great deal of structure that is not evidence of anything; every step below
exists to separate the part that is.

\subsection{Parsing, encoding, and missing data}

The file is read with \texttt{encoding='utf-8-sig'} -- a byte-order mark sits on the first
header, and without that flag the column name silently fails to match. Item codes are
parsed from the headers by regular expression rather than by renaming columns, giving four
blocks of exactly 15 items. The Likert scale maps to $-2 \ldots +2$, and the sixth option,
``No Comments'' (37 occurrences), maps to \emph{missing} rather than $0$: it is a refusal
to answer, not the substantive position ``Neutral'', and collapsing the two would place 37
refusals at the exact centre of the scale.

That leaves 542 missing entries (505 blanks plus 37 refusals). How they are distributed
matters more than how many there are:

\begin{center}
\begin{tabular}{lrrrrr}
\toprule
Missing items per respondent & 0 & 1--5 & 6--12 & 13--59 & all 60 \\
\midrule
Respondents & 68 & 17 & 1 & 5 & 5 \\
\bottomrule
\end{tabular}
\end{center}

\noindent
\textbf{495 of the 542 missing cells sit in the 10 respondents with more than 12 missing
items.} Dropping those 10 leaves \textbf{86 respondents with 47 missing cells (0.9\%)},
filled with the column mean. We checked rather than assumed that nothing heavier was
needed: substituting the median moves the correlation standard deviation from 0.1363 to
0.1361 and the edge count at $|r|>0.20$ from 258 to 260, and leaves the number of
significant eigenvalues unchanged. When two simple estimators are indistinguishable in
every downstream result, a regression or Bayesian imputer adds modelling assumptions in
exchange for no measurable change.

\subsection{Row-centring and the association measure}

The grand mean across all responses is $+1.14$ on a $-2 \ldots +2$ scale: respondents
agree with almost everything. Uncorrected, this acquiescence baseline dominates the
similarity structure, because someone who agrees with everything pushes all 60 columns up
together, which is indistinguishable from genuine item association. On the uncentred
matrix, 83\% of the 1{,}770 item correlations are positive and 696 pairs exceed
$|r|>0.20$ -- a 39\% dense graph encoding mostly who is agreeable. We therefore subtract
each respondent's own mean across all 60 items from their row before correlating, leaving
only the \emph{relative} pattern of what they lean toward given their own baseline. After
centring, 45\% of pairs are positive and 258 survive $|r|>0.20$.

Edges are Pearson correlations on that matrix \textbf{with signs retained}; we never take
absolute values. A negative edge says two questions divide the class in opposite
directions, which is a stronger claim than ``these are related'', and retaining signs is
what makes the structural-balance analysis possible at all.

\subsection{The noise floor and spectral filtering}
\label{subsec:mp}

Each of 1{,}770 item correlations is estimated from 86 respondents, so unrelated questions
scatter around zero with standard deviation $1/\sqrt{n-3} = 0.110$; across 1{,}770 pairs
many land high by luck. We quantify this by re-running the whole pipeline on data with the
association destroyed -- every column permuted independently, then centred and correlated
exactly as the real data, 200 times. Observed correlations have SD 0.136 against 0.108 for
the null (Figure~\ref{fig:noise_floor}):

\begin{center}
\begin{tabular}{lrrrr}
\toprule
Threshold $|r| >$ & 0.15 & 0.20 & 0.25 & 0.30 \\
\midrule
Observed edges & 467 & 258 & 111 & 49 \\
Expected by chance & 308 & 119 & 37 & 9 \\
Chance fraction & 66\% & 46\% & 33\% & 19\% \\
\bottomrule
\end{tabular}
\end{center}

PCA alone cannot settle this: PC1 explains only 8.3\% of variance, which distinguishes
weak structure from none not at all. Marchenko--Pastur does. For a correlation matrix from
$T$ observations of $N$ pure-noise variables the eigenvalues fall in an analytic band; at
$T=86$, $N=60$, $Q=1.433$ and $\lambda_{+}=(1+1/\sqrt{Q})^2=3.368$. \textbf{Exactly three
eigenvalues exceed it} -- 4.05, 3.98, 3.53, carrying 19.3\% of variance
(Figure~\ref{fig:mp_spectrum}) -- and the matrix is rebuilt from those three alone with the
diagonal restored to unity. Reading the loadings, they correspond to (i)~systemic
responsibility against technological optimism and traditional instruction,
(ii)~regulatory guardrails against individual autonomy, and (iii)~AI-enabled education and
industry against environmental stewardship. The third is a clean split; the first two are
mixed, and we offer them as readings rather than established constructs. One caveat
belongs with the headline: the fourth eigenvalue is 3.354 against a band edge of 3.368, so
``exactly three'' is correct but holds by 0.4\%.

\subsection{Nodes, edges, and threshold selection}

\textbf{Nodes are the 60 questions; edges are signed correlations above a threshold,
weighted by default and binarised only where a metric requires it} (path lengths, triad
counts). Questions were chosen over respondents deliberately: a respondent network must be
built by a neighbourhood rule such as $k$-nearest-neighbours, and that rule manufactures
its own communities, as Section~\ref{subsec:nulls} shows. An item network is thresholded
directly on a measured quantity and every node carries an interpretable label. Weights are
kept because correlation magnitude is measured information that binarising discards.

The threshold comes from the sweep in Figure~\ref{fig:threshold_sweep}, which rebuilds the
graph at all 46 cutoffs from 0.05 to 0.50. The rule is \emph{the highest cutoff at which
the giant component still spans more than two-thirds of the items and mean degree stays
above 4}, so community and path metrics remain meaningful; since spectral denoising has
already removed the sampling noise the chance table measures, raising the cutoff further
would penalise the same noise twice. This gives $\mathbf{|r|>0.20}$: 124 edges (60
positive, 64 negative), mean degree 4.13, a giant component of 41 of 60 items (68\%), and
modularity 0.405. Two independent controls agree there: the permutation null implies the
raw graph's 258 edges contain about 119 chance edges, leaving roughly 139 real; denoising,
which never sees the null, retains 124 -- an 11\% difference.

Panel~(d) shows why the sweep is necessary rather than decorative. Modularity climbs
steadily to 0.63 as the threshold rises, but by then the graph is 51 components and 47
isolated items. Quoting a single high modularity without the sweep would report
fragmentation as community structure.

\subsection{Null models}
\label{subsec:nulls}

A correlation-derived network has transitivity built into its construction: if A
correlates with B and B with C, A and C are pulled toward correlation automatically. That
alone manufactures clustering and modularity in data with no association whatsoever.
Degree-preserving rewiring and Erd\H{o}s--R\'enyi nulls both take the \emph{finished
graph} as given, so neither can detect structure the construction itself created. Only a
permutation null -- destroying association in the raw responses, then regenerating the
entire chain of centre, correlate, threshold and measure -- tests the claim being made.

This decides the result rather than decorating it. On the respondent network (86 nodes,
286 edges, one component, symmetrised $k$-NN at $k=4$) the observed modularity is 0.416
and average clustering 0.202:

\begin{center}
\begin{tabular}{lrrrl}
\toprule
Null model & Mean & SD & $z$ & Verdict \\
\midrule
Degree-preserving rewiring & 0.329 & 0.011 & 7.74 & passes comfortably \\
Erd\H{o}s--R\'enyi $G(n,m)$ & 0.332 & 0.011 & 7.79 & passes comfortably \\
Column permutation (full pipeline) & 0.400 & 0.019 & 0.79 & \textbf{fails} ($p=0.20$) \\
\bottomrule
\end{tabular}
\end{center}

\noindent
Against the graph-level nulls the community structure looks overwhelming; against the null
that reproduces the construction it disappears. We report this as a finding, not a
failure: those communities are an artefact of correlating and then applying a
neighbourhood rule, which is exactly why the item network is our primary object.

\subsection{Methods tried and rejected}
\label{subsec:rejected}

\begin{center}
\begin{tabular}{p{0.26\linewidth}p{0.40\linewidth}p{0.26\linewidth}}
\toprule
Approach & Measured outcome & Rejected because \\
\midrule
Mutual $k$-NN ($k=4$)
& 58 edges, 38 components, giant component 39.5\%
& Shatters the graph. \\
Raw (uncentred) similarity
& 696 of 1{,}770 pairs above $|r|>0.20$; 39\% density; 83\% positive
& Measures acquiescence, not opinion. \\
Percentile thresholding (p90)
& Items: 177 edges, 1 component. Respondents: 366 edges, 22 components, 20 isolates
& Fixes edge count by construction, so cannot say whether any edge beats chance. \\
Ordinal-regression imputation
& 47.6\% held-out accuracy, MAE 0.66
& Unjustifiable for 47 cells (0.9\%) where mean and median already agree. \\
Keeping 91 respondents rather than 86
& $Q=1.517$, $\lambda_{+}=3.283$, \emph{four} dimensions carrying 25.6\%
& Fourth dimension rests largely on imputed values. \\
\bottomrule
\end{tabular}
\end{center}

\noindent
The last row is a genuine sensitivity result and we report it rather than suppress it: the
``exactly three dimensions'' headline depends on the respondent filter. Retaining the five
respondents who answered between one and 47 items would mean imputing 242 cells with a
model measured at 47.6\% accuracy, and the extra dimension rests disproportionately on
those filled values. The frozen filter is the more defensible choice, but the finding is
conditional on it and a reader should know that.
